% ------------------------------------------------------------
\escenario{ESC-06}{Alguien captura el tráfico en una red compartida}

\begin{tabla}{|L{3.2cm}|Y|}
\fichacab
\bloque{Trazabilidad}
\parte{Característica}{QA-03 Seguridad: confidencialidad y autenticidad.}
\parte{Stakeholders}{STK-15 Especialista en ciberseguridad, STK-02 Padre o tutor.}
\parte{Concerns}{CRN-22: \emph{``El transporte es TCP en claro. Si no ciframos, cualquiera en la misma red puede leer las credenciales y secuestrar la sesión de un jugador.''} CRN-05 (datos del menor).}
\parte{Restricciones y dependencias}{CON-02 (TCP), CON-03 (sin WebSockets), CON-08 (segundo factor), DEP-02 (proveedor del segundo factor).}
\parte{Tensiones y preguntas}{TEN-01 Latencia contra integridad. Q-09 (cifrado del tráfico y mecanismo).}
\parte{Tipo y prioridad}{Exploratorio, en tiempo de ejecución (bajo ataque). Prioridad (A, A).}
\bloque{Escenario}
\parte{Fuente del estímulo}{Persona conectada a la misma red Wi-Fi que el jugador, con un analizador de tráfico.}
\parte{Estímulo}{Captura todo el tráfico entre el cliente y el servidor mientras el jugador se registra, inicia sesión con segundo factor y juega una partida con chat. Después intenta reenviar lo capturado para abrir la sesión del jugador o hacer jugadas en su nombre.}
\parte{Ambiente}{Red compartida, como la de una escuela o un café, con el juego en operación normal.}
\parte{Artefacto}{Capa de red sobre TCP y módulo de cuentas.}
\parte{Respuesta}{Lo capturado no deja leer la contraseña, el código del segundo factor, el identificador de sesión ni el contenido del chat, y no sirve para entrar ni jugar haciéndose pasar por el jugador.}
\parte{Medida de la respuesta}{En la captura de una sesión completa, buscando cadenas conocidas sembradas en la prueba, aparecen 0 contraseñas, 0 códigos de segundo factor, 0 identificadores de sesión y 0 mensajes de chat legibles. El servidor rechaza el 100~\% de 50 intentos de reenviar lo capturado. Con la protección activa, ESC-01 se sigue cumpliendo.}
\end{tabla}

\textbf{Por qué es relevante.} Con TCP directo y sin WebSockets, el juego no hereda ningún cifrado de la plataforma: si el equipo no lo resuelve, todo viaja en claro. En ese caso, el segundo factor que exige CON-08 no protege nada, porque el código se puede leer y el identificador de sesión se puede robar después de haberlo validado (Q-09 lo señala). Los usuarios son niños que juegan en redes de escuela, justo donde este ataque es más fácil, y el chat contiene datos personales de menores (CRN-05). La última medida hace visible la tensión con el desempeño: la protección cuesta tiempo en cada mensaje y no puede romper ESC-01.

\textbf{Cómo se verifica.} Se captura el tráfico de una sesión de prueba con un analizador de red y se buscan en la captura la contraseña, el código y los mensajes que se usaron durante la prueba. Después se reenvían los paquetes capturados contra el servidor y se comprueba en su registro que ninguno obtuvo sesión ni cambió una partida. Por último, se repite ESC-01 con la protección activa.
